% !PRINTED 647
\weberpage{647}
%<<CONT
But $r^s \pm r^{s'} = r^s (1 \pm r^{h-2s}) = r^s (1 \mp r^{\mu+h-2s})$ is always divisible by $1 - r$, and from this it would follow that $\mathfrak{E}(r)$ would be divisible by $1 - r$, which is not a unit, and this contradicts the concept of the unit. Hence the assumption that in formula (4) the exponent $h$ is odd is inadmissible, and, as above, we can set
\[
\mathfrak{E}(r) = r^{\tfrac12 h}\, e(r)
\]
where $e(r)$ is a real unit. With this, Theorem 1 is proved.

We wish to single out a system of real units: According to §. 169 the $\mu$ quantities $1 - r^n$ are all associated with one another. Accordingly the quotient
\[
\frac{1 - r^n}{1 - r^{n'}},
\]
in which $n, n'$ denote two numbers relatively prime to $m$, is a unit. But from this one derives, when
\[
r = e^{\frac{2\pi i}{m}}
\]
is set, the real unit
\weq{6}{\dfrac{\sin \dfrac{n\pi}{m}}{\sin \dfrac{n'\pi}{m}} = \dfrac{r^{\frac{n}{2}} - r^{-\frac{n}{2}}}{r^{\frac{n'}{2}} - r^{-\frac{n'}{2}}} = r^{\frac{n'-n}{2}}\, \dfrac{1 - r^n}{1 - r^{n'}}}
If $m$ is even, one can set $n' = n + \dfrac{m}{2}$, and obtains for this case the real units
\weq{7}{\tau_n = \mathrm{tang}\,\dfrac{n\pi}{m},}
in which $n$ can denote any odd number.

If one replaces $n$ by $\tfrac{1}{2} m - n$, then $\tau_n$ goes over into the reciprocal value. If one lets $n$ run through the sequence of numbers
\[
1,\ 3,\ 5,\ \cdots,\ \frac{m}{4} - 1
\]
then $\tau_n$ takes nothing but positive properly fractional values.

\begin{center}
\rule{0.2\linewidth}{0.4pt}
\end{center}
